\newcommand{\numgrow}{V}
\newcommand{\numgrowl}[1][\l]{\numgrow_{#1}}
\newcommand{\interl}[1][\l]{\T_{#1}}

\section{Numerical Result}
\label{sec:numerical}

Let $\l^* > 0$ be an integer and let $N = 2^{\l^* - 1} \cdot n$. Our upper bound of $\omega(1, \kappa, 1)$ is formed by successively applying \cref{thm:global-stage-with-eps,thm:consituent_MM_terms,thm:constituent-stage-with-eps} to degenerate $2^{o(n)}$ independent copies of $\CW_q^{\otimes N} \equiv \bigbk{\CW_q^{\otimes 2^{\l^* - 1}}}^{\otimes n}$ into independent matrix multiplication tensors of the form $\angbk{a, a^{\kappa}, a}$, shown in \cref{alg:framework}.

\begin{figure}[ht]
  \centering
  \begin{tcolorbox}
    \captionof{algocf}{Procedure of Degeneration}{\label{alg:framework}}
    Let $\eps > 0$ be a fixed constant and $\l^* > 0$ be an integer.
    \begin{enumerate}
    \item Degenerate $2^{o(n)}$ independent copies of $\bigbk{\CW_q^{\otimes 2^{\l^* - 1}}}^{\otimes n}$ into $\numgrowl[\l^*]$ (independent) copies of a level-$\l^*$ $(\eps \cdot 3^{\l^*})$-interface tensor $\interl[\l^*]$, where the number of copies $\numgrowl[\l^*]$ and the parameter list of $\interl[\l^*]$ are given in \cref{thm:global-stage-with-eps} and \cref{prop:global-stage-no-eps}.
    \item For each $\l = \l^*, \ldots, 2$:
      \begin{itemize}
      \item Degenerate \emph{every} $2^{o(n)}$ copies of the level-$\l$ $(\eps \cdot 3^{\l})$-interface tensor $\interl[\l]$ into $\numgrowl[\l - 1]$ independent copies of the tensor product of a level-$(\l - 1)$ $(\eps \cdot 3^{\l - 1})$-interface tensor $\interl[\l - 1]$ and some matrix multiplication tensor $\angbk{a_\l, b_\l, c_\l}$. Here, the number of copies $\numgrowl[\l - 1]$, the parameter list of $\interl[\l - 1]$ and the matrix multiplication size $\angbk{a_\l, b_\l, c_\l}$ are all given in \cref{thm:constituent-stage-with-eps} and \cref{prop:constituent-stage-no-eps}.
      \end{itemize}
    \item The level-$1$ $3\eps$-interface tensor $\interl[1]$ can degenerate into a matrix multiplication tensor, written $\angbk{a_1, b_1, c_1}$, according to \cref{thm:consituent_MM_terms}.
    \item So far, we have obtained $\numgrow \defeq \prod_{\l = 1}^{\l^*} \numgrowl[\l]$ copies of $\angbk{A, B, C} \equiv \bigotimes_{\l = 1}^{\l^*} \angbk{a_\l, b_\l, c_\l}$.
    \end{enumerate}
    We first let $n \to \infty$ and apply Sch{\"o}nhage's asymptotic sum inequality 
    (\cref{thm:schonhage-ineq-rect}) on the above degeneration, obtaining a bound on $\omega(1, \kappa, 1)$ which might depend on $\eps$; then, we let $\eps \to 0$, obtaining the bound $\omega(1, \kappa, 1) \le \omega'$ as long as
    \[
      \lim_{\eps \to 0} \lim_{n \to \infty} \numgrow^{1/n} \cdot \min\bigBK{A, B^{1/\kappa}, C}^{\omega'/n} \ge (q + 2)^{2^{\l^* - 1}}.
      \numberthis \label{eq:asymptotic_sum_inequality_in_algorithm}
    \]
  \end{tcolorbox}
  \vspace{-1.5em}
\end{figure}

\bigskip

Every degeneration step in \cref{alg:framework} requires a set of parameters, including the distribution $\alpha$ over constituent tensors, the proportions of tensor powers $A_1, A_2, A_3$ assigned to three regions, and others. If we are given an assignment to the parameters, we can precisely calculate
\[
  \lim_{\eps \to 0} \lim_{n \to \infty} \numgrowl^{1/n}, \quad \lim_{\eps \to 0} \lim_{n \to \infty} a_\l^{1/n}, \quad \lim_{\eps \to 0} \lim_{n \to \infty} b_\l^{1/n}, \quad \lim_{\eps \to 0} \lim_{n \to \infty} c_\l^{1/n}
\]
according to \cref{thm:global-stage-with-eps,thm:consituent_MM_terms,thm:constituent-stage-with-eps}. Plugging them into \eqref{eq:asymptotic_sum_inequality_in_algorithm} would verify the correctness of the claimed bound on $\omega(1, \kappa, 1)$.

\paragraph{Optimization strategy.}

Finding a set of parameters that lead to the best bound of $\omega(1, \kappa, 1)$ can be modeled as a constrained optimization problem:
\[
  \begin{array}{cl}
    \textup{minimize} & \qquad \omega' \\
    \textup{subject to} & \textup{all constraints in \cref{thm:global-stage-with-eps,thm:consituent_MM_terms,thm:constituent-stage-with-eps}} \\
                      & \displaystyle \lim_{\eps \to 0} \lim_{n \to \infty} \numgrow^{1/n} \cdot \min\bigBK{A, B^{1/\kappa}, C}^{\omega'/n} \ge (q + 2)^{2^{\l^* - 1}}.
  \end{array}
  \numberthis \label{eq:optimization_problem}
\]
We used \emph{sequential quadratic programming (SQP)} to solve this optimization problem, which is a well-known iterative approach for solving nonlinear constrained optimization. The software package SNOPT~\cite{SNOPT} is used for performing SQP. Like all other optimization methods for nonlinear optimization, SQP does not guarantee finding the global optimum or a specific convergence rate; the quality of the solution and the time performance both rely on the \emph{initial point} of the iterative process, which could be provided by the user.

For $\kappa = 1$, we take the parameters from \cite{LeGall32power} which Le Gall used to analyze $\CW_q^{\otimes 2^{\l^* - 1}}$ for square matrix multiplication, and transform it into a feasible solution to the optimization problem \eqref{eq:optimization_problem}, which we set as the initial point. Specifically, Le Gall's parameters consist of a distribution $\alpha$ over level-$\l^*$ constituent tensors (for the global stage) together with a split distribution $\alpha_{i,j,k}$ for every constituent tensor $T_{i,j,k}$ (for the constituent stages). We specify our parameters as follows:
\begin{itemize}
\item For every constituent tensor $T_{i,j,k}$ that appears in our interface tensors, we directly set $\alpha_{i,j,k}$ as its split distribution in every region, and let $A_1 = A_2 = A_3 = 1/3$, which means that all three regions are symmetric to each other.
\item The distribution used in our global stage is set to $\alpha$ as well. Other parameters are uniquely determined by these specified ones.
\item For every constituent tensor $T_{i,j,k}$ that contains a zero, say $i = 0$, we choose its complete split distributions $\splresX, \splresY, \splresZ$ that maximizes its size as an inner product tensor, i.e., maximizes $H(\splresY)$.
\item Other parameters are uniquely determined by the specified ones.
\end{itemize}
It is easy to see that these parameters form a feasible solution. Furthermore, these parameters actually lead to the same upper bound on $\omega$ as Le Gall's analysis. We start from this feasible solution and perform SQP to obtain an upper bound for $\omega = \omega(1, 1, 1)$.

For $\kappa \ne 1$, our strategy is to start with a solution for another $\kappa$ nearby. For example, it is natural to believe that a good solution for $\omega(1, 0.95, 1)$ is similar to that for $\omega(1, 1, 1)$. Therefore, we use our parameters for $\omega(1, 1, 1)$ as the initial point for optimizing the bound of $\omega(1, 0.95, 1)$, and proceed with SQP to obtain the bound for $\omega(1, 0.95, 1)$. Then, we can further start with our parameters for $\omega(1, 0.95, 1)$ to obtain parameters for $\omega(1, 0.90, 1)$, and so on.

\paragraph{Lagrange multipliers.} In \cref{thm:global-stage-with-eps}, we need to calculate $P_\alpha = \max_{\alpha' \in D} H(\alpha') - H(\alpha)$ where $D$ represents the set of distributions that share marginals with $\alpha$. Although this definition of $P_\alpha$ is not a closed form in terms of $\alpha$, we can let the max-entropy distribution $\alpha_{\max} \defeq \argmax_{\alpha' \in D} H(\alpha')$ be an optimizable variable, and use the method of Lagrange multipliers to ensure that $\alpha'$ has the largest entropy among $D$.

Formally, we first add linear constraints to force $\alpha_{\max}$ and $\alpha$ to have the same marginals:
\begin{align*}
  \sum_{j + k = 2^{\l^*} - i} \bigbk{\alpha_{\max}(i,j,k) - \alpha(i,j,k)} &= 0, \qquad \forall i = 0, 1, \ldots, 2^{\l^*},
                                                  \numberthis \label{eq:same_marginal_constraint_x} \\
  \sum_{i + k = 2^{\l^*} - j} \bigbk{\alpha_{\max}(i,j,k) - \alpha(i,j,k)} &= 0, \qquad \forall j = 0, 1, \ldots, 2^{\l^*},
                                                  \numberthis \label{eq:same_marginal_constraint_y} \\
  \sum_{i + j = 2^{\l^*} - k} \bigbk{\alpha_{\max}(i,j,k) - \alpha(i,j,k)} &= 0, \qquad \forall k = 0, 1, \ldots, 2^{\l^*},
                                                  \numberthis \label{eq:same_marginal_constraint_z} \\
  \sum_{i,j,k} \alpha_{\max}(i,j,k) &= 1, \numberthis \label{eq:alpha_max_sum=1} \\
  \alpha_{\max}(i,j,k) &\ge 0, \qquad \forall i+j+k = 2^{\lvl^*}. \numberthis \label{eq:alpha_max_nonnegative}
\end{align*}
Let $\lambda_X(i), \lambda_Y(j), \lambda_Z(k), \lambda_S$ ($0 \le i, j, k \le 2^{\l^*}$) be Lagrange multipliers for \eqref{eq:same_marginal_constraint_x}, \eqref{eq:same_marginal_constraint_y}, \eqref{eq:same_marginal_constraint_z}, \eqref{eq:alpha_max_sum=1} respectively, which we also treat as optimizable variables. Then the first-order optimality of $H(\alpha_{\max})$ can be written as
\[
  \lambda_X(i) + \lambda_Y(j) + \lambda_Z(k) + \lambda_S = \ln \alpha_{\max}(i,j,k) + 1, \qquad \forall i + j + k = 2^{\l^*}. \numberthis \label{eq:lagrange_constraint}
\]
(Note that any $\alpha_{\max}$ satisfying \eqref{eq:lagrange_constraint} will also satisfy strict inequalities in \eqref{eq:alpha_max_nonnegative}, thus we do not need to create Lagrange multipliers for \eqref{eq:alpha_max_nonnegative}.)
Since the entropy function $H(\cdot)$ is strictly concave, any $\alpha_{\max}$ satisfying these constraints is guaranteed to have maximum entropy. (Conversely, the true max-entropy distribution $\alpha_{\max}$ will satisfy all these requirements.) We include these Lagrange multiplier constraints \eqref{eq:lagrange_constraint} in our optimization problem \eqref{eq:optimization_problem}.\footnote{In the program, we use the exponential form of \eqref{eq:optimization_problem}: $\exp(\lambda_X(i) + \lambda_Y(j) + \lambda_Z(k) + \lambda_S - 1) = \alpha_{\max}(i,j,k)$, in order to avoid numerical issues like $\ln 0$.} Similarly, in \cref{thm:constituent-stage-with-eps}, we also introduce Lagrange multiplier constraints when we need to ensure that some distribution has maximum entropy given its marginals.

\paragraph{Smooth the landscape.} In \cref{thm:global-stage-with-eps,thm:constituent-stage-with-eps}, the intermediate variables named $E_1, E_2, E_3$ are minimums of three terms. If we calculate them according to the definition, it would create a ``spike'' (non-differentiable point) in the landscape, which is unfriendly for many optimizable methods including SQP. (SQP requires all objective and constraint functions to be twice continuously differentiable.) To address this issue, we treat $E_1, E_2, E_3$ as optimizable variables and transform the minimum into linear inequality constraints:
\[
  E = \min(x, y, z) \qquad \Rightarrow \qquad E \le x, \; E \le y, \; E \le z.
\]
Since $E$ (any of $E_1, E_2, E_3$) is positively correlated with the number of matrix multiplication tensors we produce, we do not need to worry that $E$ takes on a value smaller than $\min(x, y, z)$. The newly introduced constraints are linear and thus have smooth landscapes. We include these auxiliary optimizable parameters and constraints in the optimization problem \eqref{eq:optimization_problem}. In practice, we also observe that SQP would not work well without this type of smoothing.

\paragraph{Numerical results.}
\label{subsec:numerical}

We wrote a MATLAB~\cite{MATLAB2022} program to solve the optimization problem \eqref{eq:optimization_problem}, with the help of SNOPT~\cite{SNOPT}, a software package for solving large-scale optimization problems. By running the program for different $\kappa$, we obtained various upper bounds of $\omega(1, \kappa, 1)$, as shown in \cref{table:result}. All bounds are obtained by analyzing the fourth power\footnote{Our analysis also works for the eighth power, but it was too slow to solve the optimization problem due to the large number of parameters.} of the CW tensor with $q = 5$. Specifically, we obtained the important bounds $\omega \le 2.371552$, $\alpha \ge 0.321334$, and $\mu \le 0.527661$. The code and parameters are available at \url{https://osf.io/7wgh2/?view_only=ce1a6a66d9fc432d8f6da39a6ea4b6e4}.
